\section{Comparison on the selected family}\label{sz:comparison}\label{sec:comparison}

\paragraph{The sign of the common height.}
For $0\le y\le a$, the summand in $S(y,\tau)$ vanishes for $\nu>x_1$,
since $\kappa\ge a/x_1$. Every other term is bounded above by
$2y\Lambda(\nu)/\nu$. For every positive integer $h$,
$h/\nu\le2\lfloor h/\nu\rfloor$ when $\nu\le h$. Consequently
\[
 \sum_{\nu\le h}\frac{\Lambda(\nu)}{\nu}
 \le\frac{2}{h}\sum_{\nu\le h}\Lambda(\nu)\lfloor h/\nu\rfloor
 =\frac{2}{h}\log(h!)\le 2\log h.
\]
The equality counts the prime-power divisors of each integer up to $h$.
Thus $0\le S(y,\tau)\le4y\log x_1\le yL$.
Here is a direct lower bound on the archimedean term which suffices for
the cancellation. From $\mathrm{rd}_E\ge1$,
$\log\Omega\le\log(2\pi)+|c_K|\le 2+C_P\le L$ by the first scale
inequality. Since $M\ge b^{20}$, $\log(M/\Omega)\ge19L$.
For an integer $0\le y\le a$, $j(y)\ge19yL$.
Apply this with $y=\kappa(\tau)$ and sum \eqref{sz:height}.
Both $m$ and $k$ are positive by the rank and scale bounds. In particular,
\begin{equation}\label{sz:nonneg}
 H\ge0,\qquad U>0,\quad V>0 .
\end{equation}
This proof needs no estimate on the sign of $\Tot_D$.

\begin{chatAtheorem}[Signed determinant comparison]\label{sz:conclusion}
For an admissible $D$ assume the construction hypotheses of
Theorems A, B and C as specified in
Sections~\ref{A:section}, \ref{B:frontier} and \ref{c3:frontier}.
Then the signed excess of this datum satisfies \eqref{eq:det}.
\end{chatAtheorem}
\begin{proof}
Use the set $t$ and the actual ordered family from \eqref{B:witnesses}.
The vector $y_i=P_{g_i}\phi_i$ has $g_i\in t$ and
$\phi_i\in\mathcal S\subset\mathcal I$, where $\mathcal S$ is the short
set of Theorem~\ref{A:main}. Write $h=H-m\Tot_D+V$.
For a fixed $\sigma$, the matrix in \eqref{sz:covol} is the Gram matrix of
the columns $G_{\rm cl}^{1/2}v_\sigma(y_i)$ in an ordinary Euclidean
complex vector space. Gram--Schmidt gives a determinant equal to the
product of their squared orthogonal residual lengths, hence at most
$\prod_i N_\sigma(y_i)\le e^{nh}$. A dependent family has determinant
zero and satisfies the same inequality. This works for every $\sigma$
even when $h<0$. It proves $\operatorname{CovolumeLe}(y,c)$ with
\begin{equation}\label{sz:upper}
 c=n(H-m\Tot_D+V).
\end{equation}
For the index, the same $\phi_i$ are integral, with the same witnesses
$g_i\in t$. The index estimate \eqref{eq:short-index} therefore gives the budget $2|t|d_SU$.
The bound on $|t|d_S$ in \eqref{B:theorem}, multiplied by $2U>0$, gives
\begin{equation}\label{sz:budget}
 2|t|d_SU\le d_TJ,\qquad J=10U\log(1+2b^2).
\end{equation}
Increasing the budget in \eqref{sz:index} preserves the local matrices,
so $\operatorname{IndexLe}(y,\Delta)$ holds with $\Delta=d_TJ$.

The vectors are independent and $(1-4/b^2)d_T\le n\le d_T$, exactly
the rank assumptions of Theorem C. Its completed basis retains these
original first $n$ columns. Substituting \eqref{sz:upper} and
\eqref{sz:budget} into \eqref{c3:theorem} gives
\[
 d_T(H-3V)\le n(H-m\Tot_D+V)+2d_TJ.
\]
The coefficient of $H$ after rearrangement is $n-d_T\le0$, giving
\begin{equation}\label{sz:signed}
 nm\Tot_D\le(n-d_T)H+(n+3d_T)V+2d_TJ\le d_T(4V+2J).
\end{equation}
Admissibility gives $L\ge128$ and $b\ge3$, whence $4/b^2\le1/2$.
Theorem B gives $d_T>0$, so $0<d_T/2\le n\le d_T\le2n$.
Since $4V+2J\ge0$, division by $n$ bounds $m\Tot_D$ by $2(4V+2J)$.
Finally $1+2b^2\le b^3$ for $b\ge3$, so
$J\le30UL=30V$. It follows that
\[
 m\Tot_D\le2(4+60)V=128faL^4\le faL^5 .
\]
\end{proof}

\begin{corollary}[From the construction estimates to Dirichlet zeros]\label{cor:construction-grh}
Suppose that for each permissible pair $(q_0,d)$ and each $0<\eta<1/100$,
every sufficiently large integer $x$ admits at least one datum from
the construction of Theorem~\ref{thm:p1-data} at $(q_0,d,\eta,x)$
for which the hypotheses of Theorems A, B and C hold simultaneously.
Then every Dirichlet $L$-function is nonzero in $\Re s>1/2$, with the
principal pole at $1$ excepted, and its zeros in $0<\Re s<1$ lie on
$\Re s=1/2$. In particular, every nontrivial zero of $\zeta(s)$ has real
part $1/2$.
\end{corollary}
\begin{proof}
Fix $(q_0,d)$ and $\eta$. At each integer $x$ above the threshold for these
parameters choose one such datum $D$. Theorem~\ref{sz:conclusion} gives
\eqref{eq:det} on this very datum, with the signed sum in
\eqref{eq:total} and the scales of \eqref{eq:p1-data-scale}.
This verifies Hypothesis~\ref{hyp:scale}. Theorem~\ref{thm:main} gives
the stated nonvanishing and the assertion about zeros.
\end{proof}

The pointwise comparison also applies to admissible data outside the
constructed families whenever their construction hypotheses hold.
\subsection{Precision needed for completing the determinant family}\label{sec:precision}
The cost of completing the determinant family depends on the squared norm
of each added column. Keeping the bound from \S\ref{c3:auxiliary}
before its final enlargement gives the following sufficient accuracy for
the two Hilbert--Schmidt estimates.

\begin{proposition}\label{prop:weakerprecision}
Fix the permissible parameters and $0<\eta<1/100$. Work at the scales
of Theorem~\ref{thm:p1-data} for sufficiently large integer $x$.
Suppose $d_X>0$. Retain the shortness and index conclusions of Theorem A,
the restriction bound \eqref{B:codim}, and the pairing and chain-height
estimates used in Theorem C. Assume the auxiliary spanning estimates
of \S\ref{c3:auxiliary}, including the bound $x^{99}+x^{202}b^2$
there for the logarithm of the conjugate squared norm after projection.
Set
\[
                      e=\frac{1}{32b x^{202}} .
\]
In \eqref{B:near} replace $5b^{-4}$ by $e$, and in
\eqref{B:moment} replace $9b^{-4}$ by $e$. If both inequalities hold
for the same normalized map, then $m\Tot_D\le fa(\log b)^5$.
\end{proposition}
\begin{proof}
The projected lower bound with $Z=X$ gives $d_T>0$. If $F$ is a
subspace of the good target and $\alpha=\dim F/d_T$, the trace
average in \S\ref{B:one} gives some translate with projected trace
at most $d_X\sqrt{\alpha(1+e)}\le d_X(1+\alpha+e)/2$.
Applying the lower estimate to the inverse image of $F$ loses
at most $ed_X$, so \eqref{B:spread} holds with $\epsilon=e$.
The greedy argument of \S\ref{B:greedy} stops at codimension at
most $16e d_T$ and uses a set $t$ with
\[
 |t|d_X\le \tfrac52\bigl(1+\log(1/(16e))\bigr)d_T
          \le5\log(1+2b^2)d_T .
\]
Indeed $x^{202}\le b$ eventually because $\log b>x^\eta$.
Restricting to $S$ costs at most $5b^{-8}\log(1+2b^2)d_T$
more dimensions by \eqref{B:codim}. This cost is at most
$d_T/(2b x^{202})$ for large $b$, since $x^{202}\le b$.
The short-vector basis extraction therefore produces an independent
family of size $n$ with $n'=d_T-n\le d_T/(b x^{202})$.

Apply the block determinant argument of \S\ref{c3:compare} to this
family before its last rank estimate. The original columns contribute
$c+2\Delta$. Each added column costs at most $x^{99}+x^{202}b^2$
at infinity by \S\ref{c3:auxiliary} and at most $b$ at the finite
places by \eqref{c3:gammaheight}. Together with \eqref{c3:chainheight}
the height inequality reads
\[
 0\le c+2\Delta+n'(x^{99}+x^{202}b^2+b)+d_T(b+V-H).
\]
The cost of the added columns divided by $d_T$ is at most
$b+x^{-103}/b+x^{-202}\le V$ eventually, since
$V=f x b(\log b)^4$. Also $b\le V$, so $d_T(H-3V)\le c+2\Delta$.
The cardinality bound for $t$ gives the index budget \eqref{sz:budget},
and $n\ge d_T/2$. Substitution of \eqref{sz:upper} and the sign
$H\ge0$ prove \eqref{eq:det} by \eqref{sz:signed}.
\end{proof}

At fixed $\eta$, the bounds $\log b>x^\eta$ and
$x_1=\lfloor\exp(x^{\eta/6})\rfloor$ give $x^{202}=o(x_1^2)$.
Thus the accuracy $(32b x_1^2)^{-1}$ also suffices. More generally,
put $L=\log b$. A normalized error $O(b^{-1}L^{-A})$ in each
estimate suffices for any $A>202/\eta$, since $x^{202}<L^{202/\eta}$.
In particular $O(b^{-1-\delta})$ suffices for every fixed $\delta>0$.
These implications allow constants depending on the fixed parameters
and exponents, with a common bound over the selected data as $x$ varies.

